% Fragment: analytic comparison, unitary identification, and full theories.
% No preamble.  The master document supplies theorem-like environments,
% criticalnote/reviewnote, \artifact, \Rep, \Com, \FPdim, \Vec, and \HH.

\section{From vacuum unitarity to complete unitarity}
\label{sec:analytic-unitarity}

The exact strong-generation certificates concern the actual subVOA
$U\subset B\subset H=V_{E_8}$.  In particular, they establish
$\dim U_1=6$ and $U\ne B$; they do not establish positivity on all
modules or on fusion products.  We now separate those questions.
\emph{Vacuum unitarity} is the positive invariant form inherited by the
PCT-stable subVOA $U$.  \emph{Strong unitarity} means that all simple
ordinary $U$-modules admit compatible positive invariant forms.
\emph{Complete unitarity} additionally includes positivity of the
canonical Hermitian fusion forms, giving the VOA representation category
its unitary modular tensor structure.  An abstract ribbon equivalence to
a unitary category does not, by itself, prove this last positivity statement.

We use the preceding algebraic analysis as input: regularity of $U$,
the actual restrictions
\begin{equation}
 B\mathord\downarrow_U=U\oplus X,
 \qquad
 H\mathord\downarrow_U=U\oplus X\oplus Y^{\oplus2},
 \qquad h_X=h_Y=1,
 \label{eq:analytic-branching}
\end{equation}
the twelve simple algebraic $U$-modules, and the
$\operatorname{HI}(\mathbb Z/3)$ fusion ring of the all-module category
for the original extension.  Every use of these inputs inherits the
qualifications attached to their preceding proofs.  Agreement of modular
data alone is never used as a classification of tensor categories.

\begin{criticalnote}{Analytic proof boundary}
The local applications labelled A1--A3 below are \textbf{claims pending
expert review}.  Their proposed arguments are included so that the
domains, geometric identifications, and counting steps can be checked.
The cited external theorems are used with their stated hypotheses;
the issue under review is whether our construction satisfies and matches
those hypotheses.  Thus complete unitarity, the unitary ribbon
identification, and the specifically $\HH_3$ extension remain conditional
in this manuscript.  The exact finite certificates are separate accepted
inputs and are not weakened by this distinction.
\end{criticalnote}

\subsection{The subnet and its sector count}
\label{subsec:analytic-net-count}

Write $\mathcal N_U\subset\mathcal N_H$ for the actual conformal-net
inclusion.  The lattice VOA supplies strong locality and bounded
localized vertex operators; the unitary conformal subalgebra inherits
the relevant local construction.  Once regularity is available,
\href{https://link.springer.com/article/10.1007/s00039-024-00685-8}
{Tener, Corollary~4.24}, applies to the conformal inclusion and gives
finite index.  The holomorphic lattice net is completely rational;
\href{https://arxiv.org/html/math/0102196}
{Longo, Theorem~24}, gives complete rationality of the subnet.
This finite-index argument does not require every $U$-module to be unitary.

\begin{criticalnote}{A2: actual net branching and twelve sectors}
\begin{claim}[Net-sector count; pending expert review]
\label{clm:analytic-net-count}
The restriction of the holomorphic $H$ net has precisely the three
inequivalent irreducible $\mathcal N_U$ summands corresponding to
\eqref{eq:analytic-branching}, with multiplicities $1,1,2$.
Its boundary fusion category $\mathcal F_{\mathrm{net}}$ has rank six,
untwisted invertible subcategory $\Vec_{\mathbb Z/3}$, and the
$\operatorname{HI}(\mathbb Z/3)$ fusion ring.  Consequently
\[
 \Rep(\mathcal N_U)\simeq Z(\mathcal F_{\mathrm{net}}),
 \qquad
 \#\operatorname{Irr}\Rep(\mathcal N_U)=12.
\]
\end{claim}
The passage from algebraic summands to irreducible net summands, and
the transfer of the actual automorphism and endomorphism information,
must be verified at the level of closed local operators.  A shared
branching matrix does not identify two tensor categories.
\end{criticalnote}

Here is the proposed reduction of this claim.  On finite-energy vectors
of $H$, reconstructed $U$-fields are its original restricted fields.
Orthogonal projections onto its actual simple summands should commute
with their closures.  The resulting Lagrangian algebra has endomorphism
algebra
\begin{equation}
 \operatorname{End}(A_{\mathrm{net}})
 \cong\mathbb C\oplus\mathbb C\oplus M_2(\mathbb C).
 \label{eq:analytic-end-algebra}
\end{equation}
The center-induction description then gives rank six and a
noncommutative complex fusion algebra for $\mathcal F_{\mathrm{net}}$.
The actual order-three automorphisms of $H/B$ supply three invertibles.
If all six simples were invertible, the intermediate index would be
$[\mathcal N_B:\mathcal N_U]=2$.  The connected intermediate algebra
$1+X$ would have a sign automorphism.  Reconstruction on finite energies
would turn it into an automorphism of $B$ fixing $U$, contrary to the
actual calculation $\operatorname{Aut}(B/U)=1$.  Thus the invertible
group is exactly $\mathbb Z/3$.  The pointed nine-sector $B$ net and the
actual $H/B$ extension supply its untwisted associator class.

With these positive spherical inputs, the ring and indicator arguments
apply independently on the net side.  In particular, for
$d=(3+\sqrt{13})/2$ and $s=1+d^2$, the center calculation gives dimensions
\begin{equation}
 1,\ d^2,\ \underbrace{s,s,s,s}_{4},\
 \underbrace{3d,3d,3d,3d,3d,3d}_{6}.
 \label{eq:analytic-center-dimensions}
\end{equation}
The last six twists have exponents
$\{2,5,6,7,8,11\}/13$.  This is the proposed source of the net count
twelve.  It is not obtained by assigning the VOA character transformation
matrix to the net category.  The underlying deductions are recorded in
\artifact{hi_cft_20260923_next_structure_rank46.md},
\artifact{hi_cft_20260923_next_structure_rank6_indicators.md}, and
\artifact{hi_cft_20260923_next_structure_rank6_center.md}.

\subsection{Reverse reconstruction and exhaustion}
\label{subsec:analytic-reconstruction}

\href{https://arxiv.org/html/2507.20735v3#S9}
{Henriques--Tener, Theorem~9.5 and Corollary~9.7}, reconstruct unitary
modules and identify their closed smeared fields.  Remark~9.8 extends
the construction beyond the finite-multiplicity hypothesis.  This
extension matters: finite-dimensional energy eigenspaces must not be
silently assumed for the net representation one is trying to reconstruct.

\begin{criticalnote}{A3: ordinaryness, irreducibility, and linear fullness}
\begin{claim}[Reverse module comparison; pending expert review]
\label{clm:analytic-reconstruction}
For this regular $U$, reconstruction sends each irreducible
$\mathcal N_U$ sector to an irreducible ordinary unitary $U$-module.
It preserves $L_0$ and induces a fully faithful linear $*$-functor on
finite direct sums,
\[
 \mathcal R:\Rep(\mathcal N_U)\longrightarrow\Rep^{u}(U).
\]
No tensorator or compatibility with braiding is asserted here.
\end{claim}
The projection argument below must establish common-core invariance and
commutation with closed local fields, including their adjoints.  These
are our analytic deductions from reconstruction, not a quoted
braided-equivalence theorem.
\end{criticalnote}

Let $\mathcal K$ be an irreducible net representation and let
\[
 M=\bigcup_{R>0}1_{[0,R]}(L_0)\mathcal K
\]
be the dense bounded-energy space appearing in the extended construction.
Regularity decomposes this admissible weak $U$-module algebraically into
ordinary simple modules.  Choose a nonzero summand $S$, inherit its
positive form, and let $P$ be the Hilbert projection onto
$\overline S\subset\mathcal K$.

The intended argument first uses the agreement of the two Virasoro
actions, as in Henriques--Tener Corollary~9.6.  Each energy space of $S$
is finite-dimensional, so its closure adds no vector at a fixed energy.
It follows that $P$ preserves the bounded-energy domain and maps it
into $S$.  The unitary adjoint identities, tested against $S$, then give
commutation with all modes on the common domain.  The required closure
argument is that $P$ and $1-P$ reduce the closed smeared fields and
therefore the local algebras they generate.  Irreducibility of the net
sector implies $P=1$, and the finite-dimensional energy spaces then
give $M=S$.  This proves ordinaryness and simplicity if the stated
domain assertions hold.

For fullness, an algebraic isomorphism between simple unitary modules is
unitary after multiplication by a scalar, by uniqueness of the positive
invariant form.  Its bounded extension should intertwine the closed
fields by the same argument.  Conversely, net intertwiners commute with
the reconstructed modes.  Semisimplicity handles finite direct sums.
The detailed local source is
\artifact{hi_cft_20260923_next_m3_unitarity_scope.md}.

\begin{proposition}[Conditional strong unitarity]
\label{prop:analytic-strong-unitarity}
Assume Claims~\ref{clm:analytic-net-count}
and~\ref{clm:analytic-reconstruction}, and the preceding algebraic
count of twelve simple $U$-modules.  Then every simple $U$-module is
unitarizable, and its character is the character of a net sector.
\end{proposition}
\begin{proof}
Full faithfulness makes the twelve net simples pairwise inequivalent
simple ordinary $U$-modules.  There are exactly twelve algebraic simple
modules, so this injection is a bijection.  Each image carries its
reconstructed positive form, and $L_0$ is preserved.  Finite direct sums
give the corresponding statement for ordinary modules.
\end{proof}

This counting step supplies essential surjectivity only as a linear
comparison.  Neither fusion compatibility, equality of categorical
$S$ matrices, nor positivity of VOA fusion forms follows from it.

\subsection{The annulus compatibility claim}
\label{subsec:analytic-annuli}

The next step concerns the specific geometric annulus system of
\href{https://arxiv.org/html/1810.08168#S3.SS4}
{Tener, Definition~3.18 and Corollaries~3.20--3.22}.
Let $(B,A)$ be a localized pair for an interval $I$, with incoming
factor $A$ and outgoing factor $B$.  Subscripts distinguish their
actions on the vacuum and on a reconstructed module.  For an insertion
point $z\in\operatorname{int}(B,A)$, put
\[
 X_v(B,z,A)=\overline{B_UY_U(v,z)A_U}\in\mathcal N_U(I).
\]
The bar denotes the bounded extension when it exists.  The compatibility
required by Tener's definition of a module net representation is the
specific operator identity in the following claim.

\begin{criticalnote}{A1: new annulus compatibility lemma}
\begin{claim}[Annulus compatibility; pending expert review]
\label{clm:analytic-annulus}
For the actual geometric system above, choose matching paths and
central lifts in the two annulus constructions.  If $M$ is an ordinary
unitary module reconstructed from a sector $\pi_M$, then
\begin{equation}
 \pi_M\bigl(X_v(B,z,A)\bigr)
   =\overline{B_MY_M(v,z)A_M},
 \qquad v\in U,\quad z\in\operatorname{int}(B,A).
 \label{eq:analytic-annulus-identity}
\end{equation}
The right side exists as a bounded extension.  The same matching in the
vacuum identifies the net's geometric insertion with $X_v(B,z,A)$.
\end{claim}
The points requiring expert verification are modulewise matching of
the annulus systems and central lifts, density and core properties of
the inverse-finite-energy domains, and holomorphic continuation on every
component across a possible circle seam.  Agreement of vacuum operators
alone would not prove this claim.
\end{criticalnote}

We give the proposed proof in three steps.  First, geometric localization
must precede operator comparison.  The thin-boundary parametrizations in
the selected system agree outside $I$, and its common localizer is
geometric.  Its generating operators are univalent-semigroup factors
$\exp(-tL(\rho))$ and positive-energy diffeomorphism or M\"obius
operators.  The time-ordered integration in
\href{https://arxiv.org/html/2506.08684#S5}
{Henriques--Tener, formula~(5.2)}, gives the proposed comparison with the
geometric annulus action.  Constant generators give the same exponential;
thin factors must give the same diffeomorphism action.  A common path
and central lift are essential.  Using the same localizing lift for both
members of $(B,A)$ is intended to cancel deck-transformation factors.
This matching is required on every positive-energy module, not merely
on the vacuum.

Second, let $M^{\mathrm{fin}}$ be the finite-energy vectors and define
\begin{equation}
 \mathcal D_A=\{\xi\in\mathcal K_M:A_M\xi\in M^{\mathrm{fin}}\},
 \qquad
 \mathcal D_{B^*}=\{\eta\in\mathcal K_M:B_M^*\eta\in M^{\mathrm{fin}}\}.
 \label{eq:analytic-annulus-domains}
\end{equation}
For an incoming annulus, choose a small round circle inside its smooth
inner Jordan domain, containing the origin, and fill the complementary
annulus by $C$.  The proposed factorization, with the matched central
scalar retained, is
\begin{equation}
 A_MC_M=\lambda r^{L_0},\qquad 0<r<1,\quad\lambda\ne0.
 \label{eq:analytic-round-factorization}
\end{equation}
It puts finite-energy vectors in $\operatorname{Ran}(A_M)$ and gives
$C_M M^{\mathrm{fin}}\subset\mathcal D_A$.  The dense-range property of
$C_M$ would make $\mathcal D_A$ dense; the adjoint construction treats
$\mathcal D_{B^*}$.  This is the domain issue addressed by the proposed
use of Henriques--Tener Lemma~7.9.

Third, Theorem~9.5(i), together with insertion gluing in
Proposition~6.19, supplies the incoming identity after the preceding
geometric matching.  Adjoints supply the outgoing identity.  For
$\xi\in\mathcal D_A$ and $\eta\in\mathcal D_{B^*}$, the algebraic pairing
\begin{equation}
 \left\langle Y_M(v,z)A_M\xi,B_M^*\eta\right\rangle
 \label{eq:analytic-laurent-pairing}
\end{equation}
is a finite Laurent polynomial: both vectors have finite energy support,
and the mode grading permits only finitely many powers.  The geometric
insertion pairing is holomorphic on $\operatorname{int}(B,A)$.
The proposed geometric argument shows that every connected component
meets an incoming or outgoing open part.  The identity theorem would
then extend equality over the component, including across its seam.
Equality on the two dense test domains identifies the bounded extension
with $\pi_M(X_v(B,z,A))$.  Applying the same reasoning to the vacuum
finishes the proposed matching.  These details, rather than a general
strong-unitarity implication, are the content of Claim~A1.

\begin{theorem}[Conditional complete unitarity]
\label{thm:analytic-complete-unitarity}
Assume Claims~\ref{clm:analytic-net-count},
\ref{clm:analytic-reconstruction}, and~\ref{clm:analytic-annulus},
together with the preceding algebraic inputs.  Then $U$ is completely
unitary.
\end{theorem}
\begin{proof}
Proposition~\ref{prop:analytic-strong-unitarity} unitarizes every simple
module.  Claim~\ref{clm:analytic-annulus} supplies the representation
required by
\href{https://link.springer.com/article/10.1007/s00039-024-00685-8}
{Tener, Definition~2.18}, for each such module; finite direct sums
preserve the property.  His Corollary~4.6 then supplies positivity of
the canonical fusion forms.  Strong unitarity must be established first:
the unitarizable-fusion-product hypothesis in Theorem~4.4 cannot be used
to bypass the exhaustion argument.
\end{proof}

\section{Conditional unitary categorical identification}
\label{sec:analytic-category}

\begin{theorem}[Unitary Haagerup-center identification, conditional on A1--A3]
\label{thm:analytic-category}
Under the hypotheses of Theorem~\ref{thm:analytic-complete-unitarity},
and the preceding determination of the all-module fusion ring and its
untwisted pointed subcategory, there is a unitary ribbon equivalence
\[
 \Rep(U)\simeq Z(\HH_3).
\]
For the original embedded extension $H/U$, its all-module category is
tensor equivalent to one of $\HH_2$ and $\HH_3$.
\end{theorem}
\begin{proof}
Put $\mathcal C=\Rep(U)$ and let $A$ be the algebra of the actual
unitary extension $H$.  Complete unitarity allows application of
\href{https://arxiv.org/html/1908.03484#S2.SS4}
{Gui, Theorems~2.21 and~2.28}: $A$ is a standard commutative
Q-system and its modules are unitarizable.  The unitary braiding gives
each one-sided module a compatible opposite-side action, realizing
$\mathcal F=\mathcal C_A$ in the unitary bimodule construction.  Thus
$\mathcal F$ is a $C^*$ fusion category.  Here $\mathcal C_A$ means
\emph{all} modules.  The local modules form $\mathcal C_A^0\simeq\Vec$
because the extension is holomorphic, so substituting local modules
would lose the Haagerup category.

Its already determined fusion ring is
$\operatorname{HI}(\mathbb Z/3)$.  The actual untwisted pointed
subcategory makes the obstruction $c^{(0,3)}$ vanish, while the other
obstruction is in
\[
 H^1\!\left(\mathbb Z/2,H^2(\mathbb Z/3,\mathbb C^\times)\right)=0.
\]
These are the hypotheses in
\href{https://arxiv.org/html/1609.07604}
{Izumi, Definition~2.7 and Theorems~5.9 and~9.1}, giving the two
unitary tensor classes $\HH_2,\HH_3$.  No algebra structure on
$1+\rho$ is inserted as an extra assumption.
\href{https://arxiv.org/html/1102.2631#S1}
{Grossman--Snyder, Theorem~1.4}, places these categories in the same
Morita class.  Their centers are therefore braided equivalent.
The positive spherical structures make this a ribbon equivalence, and
the Lagrangian-center theorem gives
\[
 \mathcal C\simeq Z(\mathcal F)\simeq Z(\HH_3).
\]
\href{https://arxiv.org/html/1906.09710}
{Reutter, Theorem~5}, provides the unitary braided realization of the
equivalence, preserving the associated positive ribbon structures.
\end{proof}

This uses the published unitary classification as an external input;
its classification equations have not been replayed here.  It avoids
the broader pivotal census, but it does not prove that the natural
reconstruction functor $\mathcal R$ preserves Connes fusion, the VOA
tensor product, or braiding.  That particular comparison is a distinct
statement and is not needed for the theorem just proved conditionally.

\section{Conditional full theories and a chosen \texorpdfstring{$E_8$}{E8} extension}
\label{sec:analytic-full-theory}

\begin{proposition}[Diagonal full theory; conditional character identification]
\label{prop:analytic-diagonal}
The completely rational net $\mathcal N_U$ admits its canonical unitary
local diagonal full theory.  Under Claims~\ref{clm:analytic-net-count}
and~\ref{clm:analytic-reconstruction}, and the preceding twelve-module
character determination, its partition function is
\begin{equation}
 Z_{\mathrm{diag}}(\tau,\bar\tau)
   =\sum_{i=0}^{11}\chi_i(\tau)\overline{\chi_i(\tau)},
 \qquad c_L=c_R=8.
 \label{eq:analytic-diagonal-partition}
\end{equation}
\end{proposition}
\begin{proof}
Apply the diagonal Longo--Rehren construction, in the identity-coupling
case of
\href{https://arxiv.org/html/math/9911149}
{Rehren, Theorem~1.4, Proposition~1.5, and Corollary~1.6}.
Claims A2--A3 identify all sector Hilbert spaces with the twelve ordinary
modules and preserve their Hamiltonians.  Hence their traces give
\eqref{eq:analytic-diagonal-partition}.  The annulus compatibility claim
is not needed for this net construction or this linear character matching.
\end{proof}

\begin{theorem}[Chosen $\HH_3$ extension and defects; conditional on A1--A3]
\label{thm:analytic-e8-extension}
Assume the hypotheses of Theorem~\ref{thm:analytic-category} and the
preceding strict positivity of all nonvacuum lowest conformal weights.
There exists a unitary holomorphic conformal extension $H'\supset U$
isomorphic to $V_{E_8}$ whose all-module category is $\HH_3$.
For its extension algebra $A_3\in\Rep(U)$,
\begin{equation}
 \Rep(U)_{A_3}\simeq\HH_3,
 \qquad
 {}_{A_3}\Rep(U)_{A_3}\simeq\HH_3\boxtimes\HH_3^{\mathrm{op}}.
 \label{eq:analytic-extension-defects}
\end{equation}
The second category describes algebra-bimodule defects preserving the
$U$ chiral subalgebra in the corresponding full extension theory and
contains an $\HH_3$ fusion subcategory.
\end{theorem}
\begin{proof}
Let $I_{\HH_3}$ be the right adjoint of the forgetful functor
$Z(\HH_3)\to\HH_3$.  Transport the canonical algebra
$I_{\HH_3}(\mathbf1)$ along the unitary ribbon equivalence.  By the
canonical Lagrangian-algebra construction, as in
\href{https://arxiv.org/abs/1009.2117}
{Davydov--M\"uger--Nikshych--Ostrik, Section~4}, the resulting $A_3$
is connected, commutative, separable, and twist-free, with
\[
 \FPdim(A_3)^2=\FPdim(\Rep(U)),
 \qquad \Rep(U)_{A_3}^{0}\simeq\Vec,
 \qquad \Rep(U)_{A_3}\simeq\HH_3.
\]
Trivial twist makes its constituent conformal weights integral.
Connectedness and strict positivity away from the vacuum give a
one-dimensional weight-zero space.  The extension correspondence of
\href{https://arxiv.org/html/1406.3420}
{Huang--Kirillov--Lepowsky, Theorems~3.2--3.6}, together with Gui's
unitary Q-system theorem, produces a unitary conformal VOA extension
$H'$.  It is a finite sum of ordinary $U$-modules; finite-extension
$C_2$-cofiniteness and the semisimple local-module description give
strong rationality.  The Lagrangian condition makes it holomorphic.
Its central charge is eight, so
\href{https://msp.org/pjm/2004/213-2/pjm-v213-n2-p05-p.pdf}
{Dong--Mason, Theorem~1}, identifies $H'$ with $V_{E_8}$.

For the canonical algebra, the algebra-bimodule duality gives the
second equivalence in \eqref{eq:analytic-extension-defects}; see
\href{https://arxiv.org/abs/math/0111205}
{M\"uger, Sections~4 and~8}.  Here ``op'' reverses the tensor product.
The first-factor embedding $x\mapsto x\boxtimes\mathbf1$ supplies
the stated $\HH_3$ fusion subcategory.
\end{proof}

\begin{remark}[Extension choice and preserved symmetry]
\label{rem:analytic-extension-scope}
The selection of the canonical $\HH_3$ algebra proves existence of
$H'$ with the specified all-module category.  It does not identify
$A_3$ with the original embedded $H/U$ algebra; the original boundary
category may be $\HH_2$ or $\HH_3$.  Both holomorphic VOAs being
isomorphic to $V_{E_8}$ does not identify their embeddings of $U$.
Moreover, the defects in \eqref{eq:analytic-extension-defects} preserve
$U$ and need not preserve the enlarged $E_8$ chiral algebra.
The local chiral $H'$-module category is $\Vec$.
\end{remark}

The logical separation is therefore explicit: completed finite
certificates prove the actual strong-generation and properness result;
the algebraic identification supplies its stated representation-theory
inputs; A2--A3 supply sector exhaustion and strong unitarity; A1 supplies
the additional analytic comparison needed for complete unitarity.
The unitary categorical and $\HH_3$ extension conclusions retain those
conditions until the highlighted applications are independently verified.
The synthesis sources are
\artifact{hi_cft_20260923_c8two_final_theorem.md} and
\artifact{hi_cft_20260923_c8two_certificate_audit.md}; their unqualified
analytic conclusion is deliberately recorded here with the review boundary
made visible.
